% year: תשע"ז
% type: מבחן
% semester: א
% moed: ב
% number: 3ב
% points: 12
% categories: ivt, mvt
% source: exams/מבחנים/תשעז/מועד ב תשעז.pdf

\begin{question}
(12 נקודות) האם למשוואה $\cos(x)-x=0$ יש פתרון בקטע $(-1,1)$, האם הוא יחיד? צטטו במדויק כל משפט שהנכם מסתמכים עליו.
\end{question}

\begin{hint}
הגדירו $g(x)=\cos x-x$ ובדקו את הסימנים של $g(\pm1)$. ליחידות, בדקו את סימן $g'$.
\end{hint}

\begin{solution}
נגדיר $g(x)=\cos x-x$. זו פונקציה רציפה וגזירה על כל $\R$.

\textbf{קיום.} \emph{משפט ערך הביניים (בולצאנו):} אם $g$ רציפה בקטע סגור $[a,b]$ ו-$g(a),g(b)$ בעלי סימנים מנוגדים, אז קיימת $x_0\in(a,b)$ עם $g(x_0)=0$.
כאן $g$ רציפה ב-$[-1,1]$, ו-
\[
g(-1)=\cos1+1>0,\qquad g(1)=\cos1-1<0
\]
(כי $0<\cos1<1$). לכן קיים $x_0\in(-1,1)$ עם $\cos x_0-x_0=0$.

\textbf{יחידות.} $g'(x)=-\sin x-1$. מתקיים $g'(x)\le0$ לכל $x$, ושוויון $g'(x)=0$ רק כאשר $\sin x=-1$, כלומר $x=-\frac\pi2+2\pi k$; אף נקודה כזאת אינה בקטע $(-1,1)$ (כי $\frac\pi2>1$). לכן $g'(x)<0$ לכל $x\in(-1,1)$.
לפי מסקנה ממשפט לגרנז' (\emph{אם $g'<0$ בקטע אז $g$ יורדת ממש בו}), $g$ יורדת ממש ב-$(-1,1)$ ולכן חד-חד-ערכית שם, ויש לה לכל היותר אפס אחד.
(באופן שקול: אם היו שני שורשים $x_1<x_2$, לפי \emph{משפט רול} — $g$ רציפה ב-$[x_1,x_2]$, גזירה ב-$(x_1,x_2)$ ו-$g(x_1)=g(x_2)$ — הייתה קיימת $\xi\in(x_1,x_2)$ עם $g'(\xi)=0$, בסתירה.)

\textbf{תשובה:} כן, למשוואה יש פתרון בקטע $(-1,1)$, והוא יחיד.
\end{solution}
